\section*{\centering URAIAN SINGKAT}
Fisika modern saat ini sedang berada dalam perjalanan pencarian "teori segala sesuatu" dengan melakukan penggabungan seluruh gaya fundamental dalam suatu persamaan. Relativitas umum hadir untuk menjelaskan konsep gravitasi sebagai pembaharuan dari gravitasi Newton, yang mana dianggap sebagai deskripsi salah satu gaya fundamental. Ketiga gaya lain sudah diintegrasikan lewat protokol mekanika kuantum, dan saat ini masih kesulitan untuk menyatukan gaya terakhir yaitu gravitasi. Radiasi Hawking diperkenalkan sabagai konsep pertama yang menggabungkan dua kutub tersebut, yaitu relativitas umum pada lubang hitam dengan penciptaan partikel lewat mekanika kuantum. Dalam upaya pencarian tersebut, tentunya konsep baru ini menjadi salah satu titik terang. Oleh karena itu, perlu dilakukan pembahasan lebih lanjut mengenai radiasi Hawking yang terjadi pada lubang hitam selain Schwarzschild, misalnya pada lubang hitam Kerr. Namun, beberapa solusi lubang hitam tersebut muncul dari ruang-waktu yang hampa dan statis. Oleh karena itu diperkenalkan solusi Vaidya-Kerr yang bersifat non-statis. Dengan ini diharapkan akan muncul pemahaman lebih lanjut mengenai proses keluarnya informasi lubang hitam dan dapat melakukan prediksi temperatur dari sebuah lubang hitam.